\documentclass[a4paper, serif, 10pt]{moderncv}

%% moderncv
% color and style
\moderncvcolor{black}
\moderncvstyle{banking}

%% page layout/margins
\usepackage[top=2.5cm, bottom=2.5cm, left=1.5cm, right=1.5cm]{geometry}

\nopagenumbers{}

%% Babel config for Italian language
% ref:
% - https://latex3.github.io/babel/guides/locale-italian.html
\usepackage[italian]{babel}

%% fontspec
\usepackage{fontspec}

\IfFontExistsTF{Linux Libertine}{
  \setmainfont[Ligatures={TeX, Common}, Numbers=OldStyle]{Linux Libertine}
}{}
\IfFontExistsTF{Linux Libertine O}{
  \setmainfont[Ligatures={TeX, Common}, Numbers=OldStyle]{Linux Libertine O}
}{}

%% CTeX
% CJK support, used to show CJK characters in the resume
%
% - heading=false: disable ctex heading
% - punct=kaiming: use kaiming punctuations styles for CJK
% - scheme=plain: use plain scheme, do not override \normalsize font size
% - space=auto: space settings for CJK characters
%
% CJK fonts are resolved per locale language, e.g. Noto Serif/Sans CJK JP
% for Japanese resumes, with fallback to the next region if not installed
%
% ref:
% - http://ctan.mirrorcatalogs.com/language/chinese/ctex/ctex.pdf
\usepackage[UTF8, heading=false, punct=kaiming, scheme=plain, space=auto]{ctex}

\IfFontExistsTF{Noto Serif CJK SC}{
  \setCJKmainfont{Noto Serif CJK SC}
}{}
\IfFontExistsTF{Noto Sans CJK SC}{
  \setCJKsansfont{Noto Sans CJK SC}
}{}

%% line spacing
\usepackage{setspace}
\setstretch{1.125}

%% URL styling - use normal text instead of monospace
\urlstyle{same}

%% Patch \httplink and \httpslink to handle full URLs
% The original moderncv macros always prepend http:// or https:// to URLs.
% This causes issues when the URL already has a protocol (e.g., https://example.com)
% resulting in malformed URLs like https://https://example.com.
% This patch checks if the URL already contains "://" and uses it directly if so.
% Note: We use \str_set:Nx (x-expansion) to expand macros like \@homepage first.
\ExplSyntaxOn
\str_new:N \g_yamlresume_url_str

\RenewDocumentCommand{\httpslink}{O{}m}{
  \str_set:Nx \g_yamlresume_url_str {#2}
  \str_if_in:NnTF \g_yamlresume_url_str {://}
    { \url{#2} }
    { \url{https://#2} }
}

\RenewDocumentCommand{\httplink}{O{}m}{
  \str_set:Nx \g_yamlresume_url_str {#2}
  \str_if_in:NnTF \g_yamlresume_url_str {://}
    { \url{#2} }
    { \url{http://#2} }
}
\ExplSyntaxOff

\name{Luca Moretti}{}
\title{Ingegnere del Software Senior | Architetto Cloud}
\phone[mobile]{+39 333 123 4567}
\email{luca.moretti@email.it}
\homepage{https://lucamoretti.dev}

\address{Via Roma 42, Milano, Lombardia, Italia, 20121}{}{}

\extrainfo{{\small \faGithub}\ \href{https://github.com/lucamoretti}{@lucamoretti} {} {} {} • {} {} {} 
{\small \faLinkedin}\ \href{https://www.linkedin.com/in/lucamoretti}{@lucamoretti}}

\begin{document}

\maketitle

\section{Informazioni personali}

\cvline{}{\begin{itemize}
\item Ingegnere del software con oltre 8 anni di esperienza nella progettazione e sviluppo di applicazioni web scalabili e sistemi distribuiti.
\item Specializzato in architetture cloud-native, microservizi e sviluppo full-stack con particolare attenzione alla qualità del codice e alle best practice.
\item Esperienza consolidata nella leadership tecnica di team di sviluppo, nella mentorship e nella collaborazione con stakeholder di prodotto.
\item Appassionato di innovazione, automazione e miglioramento continuo dei processi di ingegneria del software.
\end{itemize}}
\section{Formazione}

\cventry{set 2013–lug 2015}
        {Laurea magistrale, Ingegneria Informatica, Punteggio: 110/110 e lode}
        {Politecnico di Milano}
        {\url{https://www.polimi.it}}
        {}
        {\begin{itemize}
\item Percorso di eccellenza focalizzato su architetture software, sistemi distribuiti e ingegneria del software.
\item Tesi sperimentale su microservizi e orchestrazione di container in ambienti cloud.
\item Partecipazione a progetti di ricerca applicata con aziende partner.
\end{itemize}
\textbf{Corsi}: Architetture Software, Sistemi Distribuiti, Basi di Dati Avanzate, Ingegneria del Software, Sicurezza Informatica, Machine Learning}

\cventry{set 2010–lug 2013}
        {Laurea triennale, Informatica, Punteggio: 108/110}
        {Università degli Studi di Bologna}
        {\url{https://www.unibo.it}}
        {}
        {\begin{itemize}
\item Solida formazione di base in informatica, algoritmi e programmazione.
\item Progetto di laurea su applicazioni web dinamiche.
\end{itemize}
\textbf{Corsi}: Programmazione, Algoritmi e Strutture Dati, Reti di Calcolatori, Sistemi Operativi, Basi di Dati}
\section{Esperienza professionale}

\cventry{mar 2020–Presente}
        {Ingegnere del Software Senior}
        {TechNova S.p.A.}
        {\url{https://www.technova.it}}
        {}
        {\begin{itemize}
\item Guida tecnica di un team di 6 sviluppatori nella progettazione di una piattaforma SaaS per la gestione documentale.
\item Definizione dell'architettura a microservizi su Kubernetes con deployment automatizzati tramite GitOps.
\item Riduzione del 40\% del time-to-market grazie all'introduzione di pipeline CI/CD e test automatizzati.
\item Collaborazione con product manager e designer per tradurre i requisiti in soluzioni tecniche scalabili.
\item Mentorship di sviluppatori junior e definizione di standard di codice e code review.
\end{itemize}
\textbf{Parole chiave}: Microservizi, Kubernetes, CI/CD, Leadership, Architettura}

\cventry{gen 2017–feb 2020}
        {Ingegnere del Software}
        {Digitalia Labs}
        {\url{https://www.digitalialabs.com}}
        {}
        {\begin{itemize}
\item Sviluppo di applicazioni web full-stack con React, Node.js e PostgreSQL.
\item Implementazione di API RESTful e integrazione con servizi di terze parti.
\item Ottimizzazione delle performance del frontend riducendo i tempi di caricamento del 30\%.
\item Adozione di pratiche Agile e partecipazione attiva a sprint planning e retrospettive.
\end{itemize}
\textbf{Parole chiave}: React, Node.js, PostgreSQL, Agile, REST}

\cventry{set 2015–dic 2016}
        {Sviluppatore Backend}
        {WebSolutions Italia}
        {\url{https://www.websolutions.it}}
        {}
        {\begin{itemize}
\item Sviluppo di backend in Java e Spring Boot per piattaforme e-commerce.
\item Progettazione e ottimizzazione di query SQL su database Oracle.
\item Implementazione di test unitari e di integrazione con JUnit e Mockito.
\end{itemize}
\textbf{Parole chiave}: Java, Spring Boot, SQL, Oracle, Testing}
\section{Lingue}

\cvline{Italiano}{Madrelingua o bilingue}
\cvline{Inglese}{Competenza professionale completa \hfill \textbf{Parole chiave}: IELTS 8.0, Business English}
\cvline{Francese}{Competenza professionale limitata}
\section{Competenze}

\cvline{Sviluppo Backend}{Esperto \hfill \textbf{Parole chiave}: Java, Python, Node.js, Spring Boot, RESTful API, Microservizi}
\cvline{Cloud e DevOps}{Avanzato \hfill \textbf{Parole chiave}: AWS, Kubernetes, Docker, Terraform, CI/CD, GitLab CI}
\cvline{Frontend}{Avanzato \hfill \textbf{Parole chiave}: React, TypeScript, JavaScript, HTML5, CSS3, Redux}
\cvline{Database}{Avanzato \hfill \textbf{Parole chiave}: PostgreSQL, MySQL, MongoDB, Redis}
\cvline{Metodologie}{Esperto \hfill \textbf{Parole chiave}: Agile, Scrum, TDD, Code Review, Mentorship}
\section{Premi}

\cventry{nov 2021}
        {TechNova S.p.A.}
        {Premio Innovazione Tecnologica}
        {}
        {}
        {Riconoscimento per il contributo alla modernizzazione dell'architettura della piattaforma aziendale e per la riduzione dei costi operativi del 25\%.}

\cventry{giu 2015}
        {Politecnico di Milano}
        {Premio Miglior Articolo}
        {}
        {}
        {Premio assegnato alla tesi magistrale su orchestrazione di microservizi in ambienti cloud.}
\section{Certificazioni}

\cventry{mar 2022}
        {Amazon Web Services}
        {AWS Certified Solutions Architect – Professional}
        {\url{https://aws.amazon.com/certification/certified-solutions-architect-professional/}}
        {}
        {}

\cventry{set 2021}
        {Cloud Native Computing Foundation}
        {Certified Kubernetes Administrator (CKA)}
        {\url{https://www.cncf.io/certification/cka/}}
        {}
        {}

\cventry{mag 2020}
        {Scrum.org}
        {Professional Scrum Master I}
        {\url{https://www.scrum.org/professional-scrum-certifications}}
        {}
        {}
\section{Pubblicazioni}

\cventry{ott 2020}
        {Architetture a Microservizi per Applicazioni Enterprise}
        {Rivista Italiana di Ingegneria del Software}
        {\url{https://www.risi.it/articoli/microservizi-enterprise}}
        {}
        {\begin{itemize}
\item Analisi delle principali sfide nella migrazione da monoliti a microservizi.
\item Presentazione di un caso di studio reale con benefici e lezioni apprese.
\item Linee guida per l'adozione di Kubernetes e pratiche DevOps.
\end{itemize}}
\section{Progetti}

\cventry{gen 2021–Presente}
        {Applicazione open-source per la gestione di note collaborative in tempo reale.}
        {CloudNotes}
        {\url{https://github.com/lucamoretti/cloudnotes}}
        {}
        {\begin{itemize}
\item Sviluppata con React, Node.js e WebSocket per la sincronizzazione in tempo reale.
\item Deployment su AWS EKS con pipeline CI/CD automatizzata.
\item Oltre 500 stelle su GitHub e utilizzata da team di sviluppo in tutto il mondo.
\end{itemize}
\textbf{Parole chiave}: React, Node.js, WebSocket, AWS}

\cventry{mar 2019–dic 2019}
        {Dashboard per il monitoraggio delle metriche di sviluppo software.}
        {DevMetrics}
        {\url{https://github.com/lucamoretti/devmetrics}}
        {}
        {\begin{itemize}
\item Integrazione con API di GitHub e GitLab per raccogliere metriche su commit, PR e issue.
\item Visualizzazione interattiva con React e Chart.js.
\item Utilizzata internamente da Digitalia Labs per migliorare i processi Agile.
\end{itemize}
\textbf{Parole chiave}: React, Chart.js, GitHub API, Agile}
\section{Interessi}

\cvline{Tecnologia}{Intelligenza Artificiale, Open Source, Cloud Computing}
\cvline{Sport}{Corsa, Ciclismo, Nuoto}
\cvline{Musica}{Chitarra, Jazz}
\section{Volontariato}

\cventry{set 2018–Presente}
        {Mentor Tecnico}
        {Coding Girls Italia}
        {\url{https://www.codinggirlsitalia.it}}
        {}
        {\begin{itemize}
\item Mentor volontario in programmi di coding per ragazze delle scuole superiori.
\item Organizzazione di workshop su sviluppo web e carriere nel settore tecnologico.
\item Supporto alla creazione di progetti finali e orientamento professionale.
\end{itemize}}
    
\end{document}